\newif\ifglobal

%%%%%%%%%%%%%%%%%%%%%%%
%%% Compile to view %%%
%%%%%%%%%%%%%%%%%%%%%%%

%%%%%%%%%%%%%%%%%%%%%%%%%%%%%%%%%%%%%%%%%%%%%%%%%%%
%%% Readme:                                     %%%
%%% https://www.overleaf.com/read/bwdjvfjksvjy/ %%%
%%%%%%%%%%%%%%%%%%%%%%%%%%%%%%%%%%%%%%%%%%%%%%%%%%%

% >>> M?: marks a module setting number or important notice

% >>> M4: Set {./relative-path} to external files for this module <<<
\def \modulefiles {./15-DBGASSIST}
% To add an external file, use:
% (Replace '---' with actual filename)
% '\input{\modulefiles/tables/---.tex}' for tables
% '\includegraphics{\modulefiles/figures/---}' for figures
% '\subfile{\modulefiles/---}' for register files


% >>> M1: This module is in global project? <<<
% 'YES' - uncomment; 'NO' - comment out
\globaltrue

\ifglobal
    \documentclass[main\_\_EN.tex]{subfiles}
   
\else
    \documentclass[openany, 10pt]{book}

    % >>> M2: In ./00-shared/preamble-trm-module, also find
    %         settings for visibility of todo notes, watermark <<<
    \usepackage{preamble-trm-module}

    % >>> M3: Temporary labels in English,
    %         you can add your own labels there <<<
    \myexternaldocument{./00-shared/temp-labels-en}

\fi %\ifglobal


\setcounter{secnumdepth}{4}


% >>> M5: If this module needs any updates in
% './00-shared/preamble-trm-module.sty'
% (include additional package, etc.), contact your module PM
% or create the file './00-shared/changelog.tex' make the
% necessary updates yourself and document them in this file <<<


\begin{document}


% Sets language into which pre-defined bits of text are translated
\selectlanguage{english}

% Do not delete this block
\ifglobal
% Add the following front matter if
% this module is outside of global project
\else
    \InputIfFileExists{readme}{}{}
    {\let\clearpage\relax \listoftodos}
    {\let\clearpage\relax \listofdonetodos}
    \tableofcontents
    %\listoftables
    %\listoffigures
\fi


%%%%%%%%%%%%%%%%%%%%%%%%%
%%% TRM Module Begins %%%
%%%%%%%%%%%%%%%%%%%%%%%%%


%%%%%%%%%%%%%%%%%%%%%%%%%%%%%%%%%%%
%%% Debug Assistant (ASSIST_DEBUG) %%%
%%%%%%%%%%%%%%%%%%%%%%%%%%%%%%%%%%%

\hypertarget{dbgassist}{}
\chapter{Debug Assistant (ASSIST\_DEBUG)}
\label{mod:dbgassist}


%%%%%%%%%%%%%%%%
%%% Overview %%%
%%%%%%%%%%%%%%%%

\section{Overview}

Debug Assistant is an auxiliary module that features a set of functions to help locate bugs and issues during software debugging.


%%%%%%%%%%%%%%%%
%%% Features %%%
%%%%%%%%%%%%%%%%

\section{Features}
\label{sec:dbgassist-features}

The Debug Assistant module has the following features:

\begin{itemize}
    \item Stack pointer (SP) monitoring
    \item Program counter (PC) logging before the CPU resets occurs
    \item CPU debugging status logging
\end{itemize}


%%%%%%%%%%%%%%%%%%%%%%%%%%%%%%
%%% Functional Description %%%
%%%%%%%%%%%%%%%%%%%%%%%%%%%%%%

\section{Functional Description}\label{sec:dbgassist-func-descr}


%%% SP Monitoring %%%

\subsection{SP Monitoring}

The Debug Assistant module can monitor the SP so as to prevent stack overflow or erroneous push/pop. When the stack pointer exceeds the minimum or maximum thresholds, the Debug Assistant will record the PC's current value and generate an interrupt. Users can then read the recorded PC value to determine which instruction caused the out of bounds access. The minimum and maximum thresholds must be configured by software.


%%% PC Logging %%%

\subsection{PC Logging}

In some cases, software developers want to know the PC at the last CPU reset. For instance, when the program is stuck and can only be reset, the developer may want to know where the program got stuck in order to debug. The Debug Assistant module can record the PC at the last CPU reset, which can be then read for software debugging. 

%%% CPU Debugging Status Logging %%%

\subsection{CPU Debugging Status Logging}
The Debug Assistant module records the CPU debugging status by providing a set of read-only registers. Please refer to \ref{mod:riscvcpu} \textit{\nameref{mod:riscvcpu}} for more information.

%%%%%%%%%%%%%%%%%%%%%%%%%%%%%
%%% Recommended Operation %%%
%%%%%%%%%%%%%%%%%%%%%%%%%%%%%

\section{Recommended Operation}\label{sec:dbgassist-recommended-operation}


%%% SP Monitoring Configuration Process %%%

\subsection{SP Monitoring}


SP bounds check monitoring:
\begin{itemize}
    \item SP exceeds the upper bound address
    \item SP exceeds the lower bound address
\end{itemize}


The configuration process for SP monitoring is as follows:

\begin{enumerate}
    \item Configure the monitored SP threshold with \hyperref[regdesc:ASSISTDEBUGCORE0SPMINREG]{ASSIST\_DEBUG\_CORE\_0\_SP\_MIN\_REG} and\\ \hyperref[regdesc:ASSISTDEBUGCORE0SPMAXREG]{ASSIST\_DEBUG\_CORE\_0\_SP\_MAX\_REG}.
    \item Configure interrupts.
    \begin{itemize}
        \item Configure \hyperref[regdesc:ASSISTDEBUGCORE0INTRENREG]{ASSIST\_DEBUG\_CORE\_0\_INTR\_EN\_REG} to enable the interrupt of a monitoring mode. 
        \item Read \hyperref[regdesc:ASSISTDEBUGCORE0INTRRAWREG]{ASSIST\_DEBUG\_CORE\_0\_INTR\_RAW\_REG} to get the interrupt status of a monitoring mode. 
        \item Configure \hyperref[regdesc:ASSISTDEBUGCORE0INTRCLRREG]{ASSIST\_DEBUG\_CORE\_0\_INTR\_CLR\_REG} to clear their interrupts.
    \end{itemize}
    \item Configure \hyperref[regdesc:ASSISTDEBUGCORE0SPMONITORENREG]{ASSIST\_DEBUG\_CORE\_0\_SP\_MONITOR\_EN\_REG} to enable the monitoring mode(s). Various monitoring modes can be enabled at the same time.
\end{enumerate}

Read \hyperref[fielddesc:ASSISTDEBUGCORE0SPPC]{ASSIST\_DEBUG\_CORE\_0\_SP\_PC} to get the PC value when an interrupt is triggered.

The interrupt of the Debug Assistant module corresponds to the interrupt source ASSIST\_DEBUG\_INTR of the interrupt matrix. For how to map the interrupt source to the CPU interrupt, please refer to the \ref{mod:intmtrx} \textit{\nameref{mod:intmtrx}}.

%%% PC Logging Configuration Process %%%

\subsection{PC Logging Configuration Process}

The CPU sends PC value to Debug Assistant. Only when \hyperref[fielddesc:ASSISTDEBUGCORE0RCDPDEBUGEN]{ASSIST\_DEBUG\_CORE\_0\_RCD\_PDEBUGEN} is 1, the PC is valid, otherwise, it is always 0. Only when \hyperref[fielddesc:ASSISTDEBUGCORE0RCDRECORDEN]{ASSIST\_DEBUG\_CORE\_0\_RCD\_RECORDEN} is 1, \hyperref[regdesc:ASSISTDEBUGCORE0RCDPDEBUGPCREG]{ASSIST\_DEBUG\_CORE\_0\_RCD\_PDEBUGPC\_REG} samples the CPU's PC, otherwise, it keeps the original value.

The description of \hyperref[regdesc:ASSISTDEBUGCORE0RCDENREG]{ASSIST\_DEBUG\_CORE\_0\_RCD\_EN\_REG} and \hyperref[regdesc:ASSISTDEBUGCORE0RCDPDEBUGPCREG]{ASSIST\_DEBUG\_CORE\_0\_RCD\_PDEBUGPC\_REG} can be found in section \ref{regdesc:ASSISTDEBUGCORE0RCDENREG} and \ref{regdesc:ASSISTDEBUGCORE0RCDPDEBUGPCREG}.

When the CPU resets, \hyperref[regdesc:ASSISTDEBUGCORE0RCDENREG]{ASSIST\_DEBUG\_CORE\_0\_RCD\_EN\_REG} will reset, while \\ \hyperref[regdesc:ASSISTDEBUGCORE0RCDPDEBUGPCREG]{ASSIST\_DEBUG\_CORE\_0\_RCD\_PDEBUGPC\_REG} will not. Therefore, the latter will keep the PC value at the CPU reset. \hyperref[regdesc:ASSISTDEBUGCORE0RCDPDEBUGSPREG]{ASSIST\_DEBUG\_CORE\_0\_RCD\_PDEBUGSP\_REG} records the SP value at the reset.


% >>> If your module has no registers, delete sample text
%       for Sections Register Summary and Registers below <<<

%%%%%%%%%%%%%%%%%%%%%%%%
%%% Register Summary %%%
%%%%%%%%%%%%%%%%%%%%%%%%

% >>> Update the sample text below and insert
%     the info on register summary into the subfile <<<

\clearpage
\hypertarget{dbgassist-reg-summ}{}
\section{Register Summary}
\label{sec:dbgassist-reg-summ}

The addresses in this section are relative to the Debug Assistant base address provided in Table \ref{tab:sysmem-base-address} in Chapter \ref{mod:sysmem} \textit{\nameref{mod:sysmem}}.

The abbreviations given in Column \textbf{Access} are explained in Section \hyperref[glossary-access-types]{\textit{Access Types for Registers}}.

\subfile{\modulefiles/15-DBGASSIST-reg-summ__EN}


%%%%%%%%%%%%%%%%%
%%% Registers %%%
%%%%%%%%%%%%%%%%%

% >>> Update the sample text below and insert
%      the info on registers into the subfile <<<

\clearpage
\section{Registers}

The addresses in this section are relative to the Debug Assistant base address provided in Table \ref{tab:sysmem-base-address} in Chapter \ref{mod:sysmem} \textit{\nameref{mod:sysmem}}.

\subfile{\modulefiles/15-DBGASSIST-reg__EN}


\end{document}
